\documentclass[10pt,a4paper]{amsart}
\usepackage{amsmath,amssymb,amsthm,graphicx}
% For commutative diagrams we use Xy-pic
\usepackage[all]{xy}

% We use 2cell for 2-commutative diagrams.
\xyoption{2cell}
% Theorem environments.
%
\theoremstyle{plain}
\newtheorem{theorem}[subsection]{Theorem}
\newtheorem{proposition}[subsection]{Proposition}
\newtheorem{lemma}[subsection]{Lemma}

\theoremstyle{definition}
\newtheorem{definition}[subsection]{Definition}
\newtheorem{example}[subsection]{Example}
\newtheorem{exercise}[subsection]{Exercise}
\newtheorem{situation}[subsection]{Situation}

\theoremstyle{remark}
\newtheorem{remark}[subsection]{Remark}
\newtheorem{remarks}[subsection]{Remarks}

\numberwithin{equation}{subsection}

% Macros
%
\def\lim{\mathop{\mathrm{lim}}\nolimits}
\def\colim{\mathop{\mathrm{colim}}\nolimits}
\def\Spec{\mathop{\mathrm{Spec}}}
\def\Hom{\mathop{\mathrm{Hom}}\nolimits}
\def\Ext{\mathop{\mathrm{Ext}}\nolimits}
\def\SheafHom{\mathop{\mathcal{H}\!\mathit{om}}\nolimits}
\def\SheafExt{\mathop{\mathcal{E}\!\mathit{xt}}\nolimits}
\def\Sch{\mathit{Sch}}
\def\Mor{\mathop{\mathrm{Mor}}\nolimits}
\def\Ob{\mathop{\mathrm{Ob}}\nolimits}
\def\Sh{\mathop{\mathit{Sh}}\nolimits}
\def\NL{\mathop{N\!L}\nolimits}
\def\CH{\mathop{\mathrm{CH}}\nolimits}
\def\proetale{{pro\text{-}\acute{e}tale}}
\def\etale{{\acute{e}tale}}
\def\QCoh{\mathit{QCoh}}
\def\Ker{\mathop{\mathrm{Ker}}}
\def\Im{\mathop{\mathrm{Im}}}
\def\Coker{\mathop{\mathrm{Coker}}}
\def\Coim{\mathop{\mathrm{Coim}}}

% Boxtimes
%
\DeclareMathSymbol{\boxtimes}{\mathbin}{AMSa}{"02}

%
% Macros for moduli stacks/spaces
%
\def\QCohstack{\mathcal{QC}\!\mathit{oh}}
\def\Cohstack{\mathcal{C}\!\mathit{oh}}
\def\Spacesstack{\mathcal{S}\!\mathit{paces}}
\def\Quotfunctor{\mathrm{Quot}}
\def\Hilbfunctor{\mathrm{Hilb}}
\def\Curvesstack{\mathcal{C}\!\mathit{urves}}
\def\Polarizedstack{\mathcal{P}\!\mathit{olarized}}
\def\Complexesstack{\mathcal{C}\!\mathit{omplexes}}
% \Pic is the operator that assigns to X its picard group, usage \Pic(X)
% \Picardstack_{X/B} denotes the Picard stack of X over B
% \Picardfunctor_{X/B} denotes the Picard functor of X over B
\def\Pic{\mathop{\mathrm{Pic}}\nolimits}
\def\Picardstack{\mathcal{P}\!\mathit{ic}}
\def\Picardfunctor{\mathrm{Pic}}
\def\Deformationcategory{\mathcal{D}\!\mathit{ef}}
\usepackage{cleveref}
\crefname{originalreference}{original reference}{original references}
\Crefname{originalreference}{Original reference}{Original references}
\makeatletter
\newlabel{lemma-compare-injectives}{{0758}{source}{}{sourcepointer-1}{}}
\newlabel{lemma-compare-injectives@cref}{{[originalreference][0][]0758}{[1][0][]source}{}{}{}}
\makeatother
\crefname{originallemma}{original lemma}{original lemmas}
\Crefname{originallemma}{Original lemma}{Original lemmas}
\makeatletter
\newlabel{sites-cohomology-lemma-pullback-same-cohomology}{{0DD8}{source}{}{sourcepointer-2}{}}
\newlabel{sites-cohomology-lemma-pullback-same-cohomology@cref}{{[originallemma][0][]0DD8}{[1][0][]source}{}{}{}}
\makeatother
\title{Cohomology on the big etale site computes cohomology on the small one, with the six comparison identities}
\author{The Stacks Project Authors}
\date{stacks-project:etale-cohomology.tex:local-snapshot-20261011; GNU FDL 1.2 or later}
\begin{document}
\maketitle
\noindent\textit{Editorial note.} All delivered mathematical text is the author's own, copied from the pinned Stacks Project chapter file etale-cohomology.tex; no statement or proof step was rewritten, reordered or removed.\par\medskip
\noindent\textit{Editorial note.} The counted proof is the author's complete argument: the proof is written out in full: the K-injective computation for the abelian cases with the composite pullback identity, and the general comparison with the factorisation of the derived pullback for the module cases.\par\medskip
\noindent\textit{Editorial note.} The label lemma-compare-cohomology names the counted statement, and the proof environment immediately following it in the source is delivered verbatim as its proof.\par\medskip
\noindent\textit{Editorial note.} The counted statement is carried by the author's own lemma environment in the chapter source.\par\medskip
\noindent\textit{Editorial note.} The statement and proof are copied from the pinned chapter file; the two cited inputs are kept as external pointers with their canonical Stacks tags.\par\medskip
\noindent\textit{Allowed standard prerequisites (editorial).}\begin{itemize}
\item Injectivity and K-injectivity of the restrictions from the big to the small site. Same chapter, the lemma stating that the pullback of an injective or K-injective abelian sheaf on the big site stays injective on a small one. Supplies the computation of part (5) through a K-injective representative.
\item Comparison of cohomology along a continuous cocontinuous functor. Cohomology on Sites, the general lemma comparing cohomology on an object with cohomology on its image under such a functor. Supplies part (6) directly.
\end{itemize}\medskip
\ifdefined\hypertarget\hypertarget{sourcepointer-1}{}\fi
\noindent\textit{Source reference.} Original reference 0758: sources/stacks-etale-cohomology/pinned-original-source.tex, source lines 21102--21102. Prerequisite pointer invoked at proof line(s) 21295; it names a standard fact of the same chapter that the counted argument uses, and every step of the counted proof built on it is retained verbatim in the delivered body.\par\medskip
\ifdefined\hypertarget\hypertarget{sourcepointer-2}{}\fi
\noindent\textit{Source reference.} Original lemma 0DD8: sources/stacks-chow/tags.txt, source lines 16758--16758. Prerequisite pointer invoked at proof line(s) 21302; it names a standard fact of a sibling chapter that the counted argument uses, and every step of the counted proof built on it is retained verbatim in the delivered body.\par\medskip
\begin{lemma}
\label{lemma-compare-cohomology}
Let $f : T \to S$ be a morphism of schemes. Then
\begin{enumerate}
\item For $K$ in $D(S_\etale)$ we have
$H^n_\etale(S, \pi_S^{-1}K) = H^n(S_\etale, K)$.
\item For $K$ in $D(S_\etale, \mathcal{O}_S)$ we have
$H^n_\etale(S, L\pi_S^*K) = H^n(S_\etale, K)$.
\item For $K$ in $D(S_\etale)$ we have
$H^n_\etale(T, \pi_S^{-1}K) = H^n(T_\etale, f_{small}^{-1}K)$.
\item For $K$ in $D(S_\etale, \mathcal{O}_S)$ we have
$H^n_\etale(T, L\pi_S^*K) = H^n(T_\etale, Lf_{small}^*K)$.
\item For $M$ in $D((\Sch/S)_\etale)$ we have
$H^n_\etale(T, M) = H^n(T_\etale, i_f^{-1}M)$.
\item For $M$ in $D((\Sch/S)_\etale, \mathcal{O})$ we have
$H^n_\etale(T, M) = H^n(T_\etale, i_f^*M)$.
\end{enumerate}
\end{lemma}

\begin{proof}
To prove (5) represent $M$ by a K-injective complex of abelian sheaves
and apply Lemma \ref{lemma-compare-injectives}
and work out the definitions. Part (3) follows from
this as $i_f^{-1}\pi_S^{-1} = f_{small}^{-1}$. Part (1) is a special
case of (3).

\medskip\noindent
Part (6) follows from the very general Cohomology on Sites, Lemma
\ref{sites-cohomology-lemma-pullback-same-cohomology}. Then part
(4) follows because $Lf_{small}^* = i_f^* \circ L\pi_S^*$.
Part (2) is a special case of (4).
\end{proof}

\end{document}
